\subsection{Exact finite coefficient models}
\begin{proof}[Proof of Theorem~\ref{thm:finite-coefficient-models}]
By Corollary~\ref{cor:canonical-ideal-matrices}, present each
$G_s$ by an idempotent with entries in $\mathcal J$. Combine their
finite object lists and pad by zeros, giving idempotents $E_s$ in
one full action matrix ring $B$. For each source $j$, choose
admissible functionals representing all entries in column $j$ of
all these projectors. Their finite pointwise join $\phi_j$ is
still admissible, and all the coefficients transport to
$q_{ij}=\phi_j(l_i)$. An empty column uses the zero functional.

Admissibility and $Q$-linearity give, in the stated order,
\begin{equation}\label{eq:finite-coefficient-supports}
 q_{ij}q_{jk}
 =\phi_k\bigl(\phi_j(l_i)l_j\bigr)\leq q_{ik}.
\end{equation}
If $Q_0=(q_{ij})$, then $Q_0^2\leq Q_0$ and
$D=I\vee Q_0$ satisfies
\[
 D^2=D,\qquad DQ_0=Q_0=Q_0D.
\]
Its diagonal entries are upper idempotents. The ordinary fiber
$C=A_D(T)$ is unital, with diagonal corners $F(d_{ii})$; the
images of the global unit in those corners are its diagonal
identity entries. Since $d_{ij}l_j\leq l_i$, coefficient-to-action
maps give the asserted unital $\pi:C\to B$.

The two displayed support products show that the subgroup $N$
is an ideal. In each product with an $N$ entry, choose its
preimage in $F(q_{jk})$ or $F(q_{ij})$ and transport the product
through $F(q_{ik})$ to the target group of $C$. Naturality
identifies this with the original product, irrespective of
injectivity of the diagonal transitions. Its image is in the
canonical ideal by admissibility of the source functional.

Put $K=\ker\pi$. If an entry $x\in F(d_{ij})$ of an element of
$K$ vanishes in $\mathcal A(l_j,l_i)$, the directed-colimit zero
criterion supplies $r\geq d_{ij}$ with $rl_j\leq l_i$ and
$t_{d_{ij},r}(x)=0$. For a representative
$y\in F(q_{jk})$ of an $N$ entry, one has
\[
 rq_{jk}=\phi_k(rl_j)\leq\phi_k(l_i)=q_{ik}.
\]
The product $xy$ can therefore first transport $x$ to $F(r)$,
where it is zero, before transporting through $F(q_{ik})$ to
$F(d_{ik})$. Every summand of the matrix product vanishes, so
$KN=0$. A symmetric annihilator assertion is not needed.

By the choice of the functionals, $E_s\in\pi(N)$. Choose a lift
$a_s\in N$. Its defect $a_s^2-a_s$ lies in $K\cap N$ and
annihilates $a_s$ on the right. Thus $a_s^3=a_s^2$, and
$e_s=a_s^2$ is an idempotent in $N$ with image $E_s$. This
exact lift needs no semiperfectness or general lifting theorem.

For any $X,Y\in\pi(N)$ and $b\in B$, an entry of $b_{jk}$
has a representative at $r$ with $rl_k\leq l_j$. The supports
of its sandwich terms satisfy
\[
 rq_{km}=\phi_m(rl_k)\leq q_{jm},\qquad
 q_{ij}rq_{km}\leq q_{ij}q_{jm}\leq q_{im}.
\]
Thus $XbY\in\pi(N)$, even for arbitrary full action arrows
$b$. If $e,f\in N$ are idempotents, choose a preimage $c\in C$
of any element of $\pi(e)B\pi(f)$ and replace it by $ecf$.
This proves surjectivity of the corner map. If $z\in eCf$ is
in $K$, then $z=zf=0$ by $KN=0$, proving injectivity.
Composition and the corner identities are preserved by $\pi$.

Idempotent matrix calculus on finite sums of representables now
identifies these mixed corners with all amplitude morphisms,
with the customary opposite convention. Conversely an idempotent
in $N$ has canonical-ideal entries after applying $\pi$, so its
amplitude is finite projective and $\mathcal J$-generated. Zero
projectors handle zero amplitudes; an empty family may use the
zero matrix ring. None of these choices collapses the specified
support. This proves all assertions.
\end{proof}

\begin{example}[no fixed finite set of action objects]
\label{ex:no-fixed-finite-objects}
The finite coefficient model in Theorem~\ref{thm:finite-coefficient-models}
cannot in general be obtained by restricting to one fixed finite
set of action objects, even for one free support and semisimple
upper coefficient rings.

Let $k$ be a nonzero field and let
$M=\langle a,b\mid ba=1\rangle$ be the bicyclic monoid. Put
$e_n=a^nb^n$ for $n\geq1$, and let $Q$ be the semiring of finite
subsets of $M$, with union and setwise multiplication. Form the ring
\[
 R=k1\oplus\bigoplus_{n\geq1}kz_n,
 \qquad z_n^2=z_n,\quad z_nz_m=0\quad(n\ne m).
\]
Give it the $M$-grading $R_1=k1$, $R_{e_n}=kz_n$, with all
other components zero. Set $F(S)=\bigoplus_{m\in S}R_m$, with
inclusion transitions and the restricted multiplication of $R$,
and reconstruct $T=(Q,F)$. Every upper multiplicative idempotent
$E$ of $Q$ has $F(E)$ a finite product of copies of $k$.

Nevertheless, on the single free support $L=Q$ there are pairwise
nonisomorphic finitely generated projective $T$-semimodules with
amplitudes
\[
 G_n(S)=
 \begin{cases}k,&a^n\in S,\\0,&a^n\notin S,\end{cases}
 \qquad n\geq1,
\]
whose nonzero transitions are identities and whose amplitude
endomorphism rings are $k$. For any finite collection of objects
$S_1,\ldots,S_r\in L$, some $G_n$ restricts to the zero functor
on their full action subcategory, just as the zero-amplitude
projective on $L$ does.
\end{example}
\begin{proof}
Every element of $M$ has a normal form $a^ib^j$. The action on
$\mathbb N$ given by $a(t)=t+1$ and $b(t)=\max(t-1,0)$ sends
this form to $i+\max(t-j,0)$ and distinguishes the forms.
Consequently the degree $i-j$ is additive, and
$e_ne_m=e_{\max(n,m)}$. The displayed grading is multiplicative:
different nonidentity homogeneous components have zero product.
This verifies the coefficient system and hence the reconstruction.

An upper idempotent $E$ is a finite submonoid containing $1$.
An element of nonzero integer degree would have infinitely many
distinct powers. Thus $E=\{1\}\cup\{e_n:n\in A\}$ for a
finite set $A$ of positive integers. Conversely every such set
is a submonoid. The ring isomorphism
\[
 F(E)\longrightarrow k^{1+|A|},\qquad
 c_0 1+\sum_{n\in A}c_nz_n\longmapsto
       \bigl(c_0,(c_0+c_n)_{n\in A}\bigr)
\]
proves the assertion about upper coefficient rings.

The element $t_n=(\{e_n\},z_n)$ is multiplicatively idempotent.
Right multiplication by $t_n$ is an idempotent endomorphism of
the regular left $T$-semimodule, so $P_n=Tt_n$ is finitely
generated projective. Its support is $Q\{e_n\}$, identified
with $Q$ by the mutually inverse $Q$-linear maps
\[
 S\longmapsto S\{b^n\},\qquad
 U\longmapsto U\{a^n\}.
\]
Here $U\{e_n\}=U$, and the inverse identities follow from
$b^na^n=1$ and $a^nb^n=e_n$.

The fixed coefficients of right multiplication by $t_n$ at
such a support $U$ are $\{x\in F(U):xz_n=x\}$. Linear
independence of $1,z_1,z_2,\ldots$ makes this group $kz_n$
if $e_n\in U$, and zero otherwise. Normal forms show that
$e_n\in S\{b^n\}$ exactly when $a^n\in S$. This proves
the asserted amplitude, including its action and transition
maps inherited from the retract. Every amplitude endomorphism
is $k$-linear. The identity transitions from $\{a^n\}$ to
every nonzero fiber force a common scalar, and every scalar
commutes with the action. Thus its endomorphism ring is $k$.

There is no nonidentity $Q$-linear automorphism of the regular
support $Q$. Such an automorphism is right multiplication by
a finite set $U$ with a finite inverse $V$, so
$UV=VU=\{1\}$. For every $u\in U$ and $v\in V$ one has
$uv=vu=1$. Normal forms imply that $M$ has only the unit $1$:
if $a^ib^j a^rb^s=1$, then $i=s=0$ and $j=r$, and the
reverse product is $1$ only when $j=0$. Hence
$U=V=\{1\}$. The fibers at $\{a^n\}$ now distinguish
the $P_n$ even without a chosen support identification.

Right multiplication by $(\{1\},0)$ is the zero-amplitude
projector with support $Q$. Finally, a finite union
$S_1\cup\cdots\cup S_r$ omits some $a^n$. All values of
$G_n$ on these objects vanish, and therefore so do all its
restricted arrows. It remains nonzero at $\{a^n\}$. This
proves the restriction obstruction. It makes no assertion
against every possible universal ordinary-ring construction.
\end{proof}

